\documentclass[11pt]{article}
\usepackage{times}
\usepackage{fullpage}
\usepackage{url}
\usepackage{hyperref}
\usepackage{fancyhdr}
\usepackage{graphicx}
\usepackage{enumitem}
\usepackage{subcaption}
\usepackage{amsmath, amsfonts, amsthm, fouriernc}
\usepackage{IEEEtrantools}
\usepackage{parskip}
\usepackage{booktabs}

\pagestyle{fancy}
\addtolength{\headheight}{2em}
\addtolength{\headsep}{1.5em}
\lhead{\large{ME 2801}}
\rhead{\large{Solutions}}

\usepackage{sectsty}
\sectionfont{\fontsize{12}{8}\selectfont}

\begin{document}

\begin{center}
  \Large{\bf{Solutions: Laplace Transforms}}
\end{center}

Table references are to the Laplace pairs and properties handout
(equivalently, Nise Tables 2.1 and 2.2).  Every $f(t)$ below is understood
to be multiplied by the unit step $u(t)$.  Numerical verification of Parts~2--4
is in \texttt{assignment\_solutions.m}.

% ---------------------------------------------------------------
\section*{Part 1: Forward Laplace Transform}

\begin{enumerate}

\item $f(t) = 4e^{-3t}$.  Exponential pair with $a = 3$, scaled by 4:
\[
  \boxed{F(s) = \frac{4}{s+3}}
\]

\item $f(t) = 2\sin(5t)$.  Sine pair with $\omega = 5$:
\[
  \boxed{F(s) = \frac{2\cdot 5}{s^2 + 25} = \frac{10}{s^2 + 25}}
\]

\item $f(t) = 3\cos(2t)$.  Cosine pair with $\omega = 2$:
\[
  \boxed{F(s) = \frac{3s}{s^2 + 4}}
\]

\item $f(t) = t\,e^{-2t}$.  Ramp pair $t \leftrightarrow 1/s^2$, then the
frequency-shift property $e^{-at}f(t) \leftrightarrow F(s+a)$ with $a = 2$:
\[
  \boxed{F(s) = \frac{1}{(s+2)^2}}
\]

\item $f(t) = 5 - 5e^{-t}$.  Linearity: transform each term separately.
\[
  F(s) = \frac{5}{s} - \frac{5}{s+1}
       = \frac{5(s+1) - 5s}{s(s+1)}
  \quad\Rightarrow\quad
  \boxed{F(s) = \frac{5}{s(s+1)}}
\]
Either form is acceptable.  The combined form shows the two poles at $s=0$
and $s=-1$ that generate the constant and the decaying exponential.

\end{enumerate}

% ---------------------------------------------------------------
\section*{Part 2: Inverse Laplace Transform by Partial Fractions}

\begin{enumerate}[resume]

\item $F(s) = \dfrac{3}{s(s+3)}$.  Two distinct real poles:
\[
  \frac{3}{s(s+3)} = \frac{A}{s} + \frac{B}{s+3},
  \qquad
  A = \left.\frac{3}{s+3}\right|_{s=0} = 1,
  \qquad
  B = \left.\frac{3}{s}\right|_{s=-3} = -1.
\]
\[
  \boxed{f(t) = 1 - e^{-3t}}
\]

\item $F(s) = \dfrac{s+1}{(s+2)(s+4)}$.
\[
  \frac{s+1}{(s+2)(s+4)} = \frac{A}{s+2} + \frac{B}{s+4},
  \qquad
  A = \left.\frac{s+1}{s+4}\right|_{s=-2} = \frac{-1}{2},
  \qquad
  B = \left.\frac{s+1}{s+2}\right|_{s=-4} = \frac{-3}{-2} = \frac{3}{2}.
\]
\[
  \boxed{f(t) = -\tfrac{1}{2}e^{-2t} + \tfrac{3}{2}e^{-4t}}
\]
Check: $f(0^+) = -\tfrac12 + \tfrac32 = 1$, which matches the initial-value
theorem, $\lim_{s\to\infty} sF(s) = 1$.

\item $F(s) = \dfrac{2}{s^2 + 4s + 3}$.  Factor the denominator first:
$s^2 + 4s + 3 = (s+1)(s+3)$.
\[
  \frac{2}{(s+1)(s+3)} = \frac{A}{s+1} + \frac{B}{s+3},
  \qquad
  A = \left.\frac{2}{s+3}\right|_{s=-1} = 1,
  \qquad
  B = \left.\frac{2}{s+1}\right|_{s=-3} = -1.
\]
\[
  \boxed{f(t) = e^{-t} - e^{-3t}}
\]

\item $F(s) = \dfrac{10}{s(s^2 + 4s + 8)}$.  The quadratic has complex
roots ($s = -2 \pm 2j$), so keep it as a quadratic factor:
\[
  \frac{10}{s(s^2+4s+8)} = \frac{A}{s} + \frac{Bs + C}{s^2+4s+8}.
\]
Cover-up gives $A = 10/8 = 5/4$.  Clearing denominators,
\[
  10 = A(s^2 + 4s + 8) + (Bs + C)s,
\]
and matching coefficients: $s^2$: $A + B = 0 \Rightarrow B = -5/4$;
$s^1$: $4A + C = 0 \Rightarrow C = -5$.  So
\[
  F(s) = \frac{5/4}{s} - \frac{5}{4}\cdot\frac{s+4}{s^2+4s+8}.
\]
Complete the square, $s^2 + 4s + 8 = (s+2)^2 + 2^2$, and split the numerator
so that it matches the shifted cosine and sine pairs:
\[
  \frac{s+4}{(s+2)^2 + 2^2}
  = \frac{s+2}{(s+2)^2 + 2^2} + \frac{2}{(s+2)^2 + 2^2}
  \;\leftrightarrow\;
  e^{-2t}\cos(2t) + e^{-2t}\sin(2t).
\]
\[
  \boxed{f(t) = \frac{5}{4}\Bigl[\,1 - e^{-2t}\bigl(\cos 2t + \sin 2t\bigr)\Bigr]}
\]
Checks: $f(0^+) = \tfrac54(1 - 1) = 0$ and
$f(\infty) = \tfrac54 = \lim_{s\to 0} sF(s) = 10/8$.
Using $\cos\theta + \sin\theta = \sqrt{2}\cos(\theta - \pi/4)$, the same
result can be written $f(t) = \tfrac54\bigl[1 - \sqrt{2}\,e^{-2t}\cos(2t - \pi/4)\bigr]$.

\end{enumerate}

% ---------------------------------------------------------------
\section*{Part 3: Solving ODEs via Laplace Transforms}

Zero initial conditions, so $\mathcal{L}\{\dot c\} = sC(s)$ and
$\mathcal{L}\{\ddot c\} = s^2 C(s)$.  The input is $R(s) = 1/s$.

\begin{enumerate}[resume]

\item $\dot{c} + 4c = 8\,r$.

(a)~$sC(s) + 4C(s) = 8R(s)$.

(b)~$C(s) = \dfrac{8}{s+4}\cdot\dfrac{1}{s} = \dfrac{8}{s(s+4)}$.

(c)~$\dfrac{8}{s(s+4)} = \dfrac{2}{s} - \dfrac{2}{s+4}$, so
\[
  \boxed{c(t) = 2\bigl(1 - e^{-4t}\bigr)}
\]
The response rises to a final value of $2$ (the DC gain $8/4$) with time
constant $\tau = 1/4$\,s.

\item $\ddot{c} + 5\dot{c} + 6c = 6\,r$.

(a)~$(s^2 + 5s + 6)\,C(s) = 6R(s)$.

(b)~$C(s) = \dfrac{6}{s(s^2 + 5s + 6)} = \dfrac{6}{s(s+2)(s+3)}$.

(c)~Three distinct real poles:
\[
  \frac{6}{s(s+2)(s+3)} = \frac{A}{s} + \frac{B}{s+2} + \frac{C}{s+3},
\]
\[
  A = \left.\frac{6}{(s+2)(s+3)}\right|_{s=0} = 1,
  \quad
  B = \left.\frac{6}{s(s+3)}\right|_{s=-2} = \frac{6}{(-2)(1)} = -3,
  \quad
  C = \left.\frac{6}{s(s+2)}\right|_{s=-3} = \frac{6}{(-3)(-1)} = 2.
\]
\[
  \boxed{c(t) = 1 - 3e^{-2t} + 2e^{-3t}}
\]
Checks: $c(0) = 1 - 3 + 2 = 0$ and $\dot c(0) = 6 - 6 = 0$, consistent with
the zero initial conditions.  The response is overdamped (two real poles) and
settles at the DC gain $6/6 = 1$.

\item $\ddot{c} + 4\dot{c} + 8c = 8\,r$.

(a)~$(s^2 + 4s + 8)\,C(s) = 8R(s)$.

(b)~$C(s) = \dfrac{8}{s(s^2 + 4s + 8)}$.

(c)~This is $8/10$ times the expression in Problem~9, so by linearity
\[
  c(t) = \frac{8}{10}\cdot\frac{5}{4}\Bigl[1 - e^{-2t}(\cos 2t + \sin 2t)\Bigr]
       = 1 - e^{-2t}\bigl(\cos 2t + \sin 2t\bigr).
\]
Writing $\cos 2t + \sin 2t = \sqrt{2}\cos(2t - \pi/4)$ puts it in the
requested form $c(t) = K[1 - A\,e^{-\sigma t}\cos(\omega_d t - \phi)]$:
\[
  \boxed{c(t) = 1 - \sqrt{2}\,e^{-2t}\cos\!\bigl(2t - \tfrac{\pi}{4}\bigr)}
  \qquad
  K = 1,\; A = \sqrt{2},\; \sigma = 2,\; \omega_d = 2,\; \phi = \pi/4.
\]
The poles are $s = -2 \pm 2j$, so $\sigma$ is the real part and $\omega_d$
the imaginary part.  In standard second-order form, $\omega_n = \sqrt{8}$ and
$\zeta = 4/(2\sqrt{8}) = 1/\sqrt{2} \approx 0.707$: underdamped, with the
oscillation decaying inside the envelope $\sqrt{2}\,e^{-2t}$.

\end{enumerate}

% ---------------------------------------------------------------
\section*{Part 4: Pole Locations and Time-Domain Behavior}

\begin{enumerate}[resume]

\item Poles and step-response character for each system:

\begin{center}
\begin{tabular}{@{}llp{2.2cm}p{8.2cm}@{}}
\toprule
 & $G(s)$ & Poles & Stability / step response \\
\midrule
(a) & $\dfrac{1}{(s+1)(s+3)}$ & $-1,\;-3$ & Stable.  Overdamped: sum of two decaying exponentials, no overshoot. \\[2ex]
(b) & $\dfrac{1}{s^2 + 2s + 5}$ & $-1 \pm 2j$ & Stable.  Underdamped ($\zeta = 1/\sqrt5 \approx 0.45$): decaying oscillation with overshoot. \\[2ex]
(c) & $\dfrac{1}{s(s+2)}$ & $0,\;-2$ & Marginally stable (pole at the origin).  Step response grows without bound: after the $e^{-2t}$ transient dies out it ramps at slope $1/2$. \\[2ex]
(d) & $\dfrac{1}{(s-1)(s+3)}$ & $+1,\;-3$ & Unstable (right-half-plane pole).  Step response grows exponentially as $e^{t}$. \\
\bottomrule
\end{tabular}
\end{center}

For (c), the pole at the origin is an integrator: $C(s) = 1/\bigl(s^2(s+2)\bigr)$ has a
double pole at $s = 0$, giving $c(t) = \tfrac12 t - \tfrac14 + \tfrac14 e^{-2t}$.
The system is not asymptotically stable, but the divergence is a ramp rather than
an exponential, and there is no oscillation.

\item MATLAB.  Each system is a \texttt{tf} object; \texttt{step()} on a
stable system runs until the response settles, so for (c) and (d) pass an
explicit time vector to see the growth.
\begin{verbatim}
s = tf('s');
Ga = 1/((s+1)*(s+3));
Gb = 1/(s^2 + 2*s + 5);
Gc = 1/(s*(s+2));
Gd = 1/((s-1)*(s+3));

t = 0:0.01:6;
figure(1); clf
subplot(2,2,1); step(Ga, t); title('(a) overdamped')
subplot(2,2,2); step(Gb, t); title('(b) underdamped')
subplot(2,2,3); step(Gc, t); title('(c) pole at origin: ramps')
subplot(2,2,4); step(Gd, t); title('(d) unstable: e^t growth')

pole(Ga), pole(Gb), pole(Gc), pole(Gd)
\end{verbatim}
Expected plots: (a) a smooth monotonic rise to $1/3$; (b) an oscillatory rise
to $1/5$ with about 20\,\% overshoot; (c) a curve that becomes a straight line
of slope $1/2$; (d) a curve that bends upward without limit.  The
\texttt{pole()} output confirms the locations in the table.

\end{enumerate}

\end{document}
